We study leverage as a structural metric for software architecture. We model an architecture by its degrees of freedom (DOF), capability set, and modification surface, and define leverage as the ratio of capabilities to independent change points. In this setting we prove a software-engineering theorem arc: architectures with lower DOF have lower exact error probability under the derived reliability model, lower expected modification cost, and better leverage in fixed capability classes.

The resulting framework yields an actionable decision rule: among architectures meeting the same requirements, prefer the one with maximal leverage. We show that common software patterns fit this rule uniformly. Single-source-of-truth designs, nominal contracts, generic abstractions, convention-over-configuration, and normalization all realize the same structural mechanism: consolidating independently drifting definition loci while preserving capabilities.

We instantiate the framework on representative architectural choices and on an OpenHCS refactoring case, where scattered runtime checks collapse to a single nominal contract. The paper is accompanied by a Lean~4 artifact formalizing the leverage core, the probability model, and the main ordering theorems. This TOSEM-oriented version isolates the software-engineering contribution from the mathematical-physics results developed in the companion split.
